\section{Results}\label{sec:res}
Unless noted, $n{=}3000$, $\sigma{=}1$, $\gamma{=}1$ (GATE $=2$), error bars/shaded values are 95\% bootstrap CIs over replicates, and ``bias'' is the Monte-Carlo mean error (its standard error is reported where needed). All numbers come from the logged runs in \texttt{results/} (reproduction map in Appendix~\ref{app:repro}).

\textbf{R1: Overall accuracy (Table~\ref{tab:e2}, Fig.~\ref{fig:e2app} in the appendix).} Ignoring interference (DM) under Bernoulli designs has relative RMSE $\approx0.50$ in all 48 population$\times$mechanism cases: the bias is the entire spillover ($-\gamma$) and does not shrink with $n$. Regression estimators with an exposure feature cut this to a median of 0.16--0.32 (OLS-x) per graph under Bernoulli designs, with the best case (oracle features) at 0.09--0.20 but a worst case up to 0.95 (Table~\ref{tab:e2}). Under cluster randomisation the answer depends on the graph: on graphs with a small edge cut (WS, SBM, ego-Facebook: cut $\le0.16$) DM itself drops to 0.12--0.17 and regression gives 0.10--0.19, i.e.\ regression helps little over DM there (on ego-Facebook DM is the best of the six methods). On graphs with poor clusters (ER/BA, cut $0.82$; Stanford3, cut $0.64$) cluster randomisation brings no gain: ER and BA Louvain clusters ($K{=}110$) leave 82\% of edges cut, and Stanford3's 57 clusters give a large cluster-level variance (OLS-x SD $\approx0.5$--$0.6$ vs.\ $\approx0.15$--$0.18$ under Bernoulli, excluding \texttt{peer}).

\begin{table*}[t]
\caption{Relative RMSE (RMSE/|GATE|): median over the eight mechanisms (min--max in parentheses); $n$ as listed (3000 synthetic; 4039 egoFB; 11{,}586 Stanford3), 200 replicates per cell. Bold: lowest median per row. Bern: Bernoulli(0.5); Clust: cluster randomisation with Louvain clusters ($K$ clusters, fraction of edges cut).}\label{tab:e2}
\centering\small
\setlength{\tabcolsep}{3pt}
\begin{tabular}{llllllll}
\toprule
Graph & Design & DM & OLS-x & OLS-quad & Oracle & GBDT & H\'ajek \\
\midrule
ER & Bern & 0.50 (0.48--0.50) & 0.16 (0.13--0.27) & 0.15 (0.12--0.27) & \textbf{0.12 (0.06--0.56)} & 0.40 (0.37--0.46) & 0.33 (0.31--0.42) \\
ER & Clust ($K{=}110$, cut 0.82) & 0.42 (0.39--0.45) & 0.19 (0.16--0.27) & 0.17 (0.15--0.26) & \textbf{0.15 (0.08--0.60)} & 0.35 (0.32--0.44) & 0.28 (0.21--0.37) \\
\midrule
BA & Bern & 0.50 (0.49--0.50) & 0.19 (0.12--0.28) & 0.19 (0.12--0.28) & \textbf{0.12 (0.05--0.55)} & 0.37 (0.35--0.43) & 0.33 (0.24--0.42) \\
BA & Clust ($K{=}111$, cut 0.82) & 0.41 (0.38--0.44) & 0.19 (0.13--0.26) & 0.19 (0.13--0.26) & \textbf{0.12 (0.05--0.61)} & 0.31 (0.26--0.41) & 0.28 (0.19--0.35) \\
\midrule
WS & Bern & 0.50 (0.48--0.50) & 0.20 (0.18--0.33) & 0.20 (0.18--0.33) & \textbf{0.18 (0.08--0.36)} & 0.42 (0.36--0.61) & 0.34 (0.28--0.42) \\
WS & Clust ($K{=}49$, cut 0.16) & 0.12 (0.09--0.20) & 0.10 (0.09--0.18) & 0.10 (0.09--0.18) & \textbf{0.10 (0.08--0.16)} & 0.34 (0.32--0.57) & 0.11 (0.09--0.19) \\
\midrule
SBM & Bern & 0.50 (0.49--0.50) & 0.21 (0.17--0.38) & 0.20 (0.17--0.37) & \textbf{0.18 (0.07--0.95)} & 0.44 (0.41--0.67) & 0.33 (0.26--0.47) \\
SBM & Clust ($K{=}30$, cut 0.10) & 0.12 (0.11--0.21) & 0.11 (0.10--0.20) & \textbf{0.10 (0.10--0.20)} & 0.10 (0.10--0.20) & 0.22 (0.21--0.35) & 0.12 (0.11--0.21) \\
\midrule
egoFB & Bern & 0.50 (0.49--0.50) & 0.32 (0.23--0.52) & 0.26 (0.22--0.46) & \textbf{0.20 (0.09--0.68)} & 0.36 (0.33--0.61) & 0.36 (0.27--0.51) \\
egoFB & Clust ($K{=}27$, cut 0.07) & \textbf{0.17 (0.15--0.29)} & 0.19 (0.17--0.35) & 0.21 (0.18--0.37) & 0.20 (0.17--0.34) & 0.52 (0.50--1.08) & 0.17 (0.16--0.31) \\
\midrule
Stanford3 & Bern & 0.50 (0.50--0.50) & 0.18 (0.08--0.46) & 0.17 (0.07--0.46) & \textbf{0.09 (0.04--0.59)} & 0.23 (0.20--0.32) & 0.31 (0.19--0.50) \\
Stanford3 & Clust ($K{=}57$, cut 0.64) & 0.32 (0.27--0.39) & 0.29 (0.25--0.47) & \textbf{0.26 (0.24--0.43)} & 0.30 (0.16--0.76) & 0.27 (0.26--0.44) & 0.31 (0.25--0.44) \\
\bottomrule
\end{tabular}

\end{table*}


\begin{figure*}[t]
\centering
\includegraphics[width=\textwidth]{../figs/fig_e1.pdf}
\caption{Noise term of the OLS estimator. (a) RMSE vs mean degree on ER graphs under Bernoulli(0.5) with correctly specified linear exposure (``comm'': strength of community-level baseline effects; 300 replicates). DM is flat at $\approx1$ (bias $=-\gamma$). (b) Empirical SD of the noise term $\sum_iw_i\varepsilon_i$ (obtained by refitting on noiseless outcomes) vs.\ the prediction $\sigma\|w\|_2$ (diagonal). (c) RMSE vs.\ $n$; the dotted line has slope $-1/2$.}\label{fig:e1}
\end{figure*}

\textbf{R2: The noise term is predicted exactly; error grows with degree (Fig.~\ref{fig:e1}).} Across all 2{,}010 (configuration, OLS estimator) pairs of Experiments 1--6 and 8--10 (\texttt{fig\_diag.py}) the empirical SD of $\sum_iw_i\varepsilon_i$ divided by $\sigma\|w\|_2$ has median $1.001$ (5--95\% range $0.92$--$1.09$; log-log correlation $0.999$), as Prop.~\ref{prop:decomp} implies. With $\|w\|_2^2\approx(\bar d{+}1)/(np(1{-}p))$, RMSE of a \emph{correctly specified} regression grows from $0.10$ ($\bar d{=}4$) to $0.56$ ($\bar d{=}128$) at $n{=}3000$ (OLS-x, no community effect) (Fig.~\ref{fig:e1}a), and falls like $n^{-1/2}$ (Fig.~\ref{fig:e1}c); on these (irregular, Poisson-degree) ER graphs $n\|w\|_2^2$ divided by $(\bar d{+}1)/(p(1{-}p))$ is $0.79$ at $\bar d{=}4$, $0.95$ at $\bar d{=}16$, $1.04$ at $\bar d{=}128$, so Prop.~\ref{prop:bern} is accurate for moderate/large degree and optimistic for sparse graphs with heterogeneous degrees; $\|w\|_1$ grows like $\sqrt d$ and is constant in $n$ ($6.4$ at $\bar d{=}16$ for every $n$ from 500 to 16{,}000). The RMSE is larger than the noise term alone, because outcome heterogeneity that the regression does not model (unit baselines and community effects) is projected through the same weights: with community effects the assignment-induced SD (``sdS'') exceeds the pure-noise SD (e.g.\ $0.22$ vs.\ $0.14$ at $\bar d{=}16$). Consequently RMSE is barely sensitive to $\sigma$ in our baseline (Appendix Fig.~\ref{fig:e5}): doubling $\sigma$ from 1 to 2 changes OLS-x RMSE from $0.39$ to $0.44$ (linear case).

\begin{figure*}[t]
\centering
\includegraphics[width=\textwidth]{../figs/fig_e9.pdf}
\caption{Bias and RMSE vs.\ Bernoulli treatment probability $p$ on the SBM graph ($n{=}3000$, $\bar d{=}16$). Dotted: prediction of the linear-in-$x$ regression from Prop.~\ref{prop:bern} (finite-$d$ binomial version), no fitted parameters. 200 replicates per point.}\label{fig:e9}
\end{figure*}

\textbf{R3: Bias is the slope-at-$p$ error and does not vanish with $n$ (Fig.~\ref{fig:e9}, \ref{fig:e6}; Prop.~\ref{prop:bern}).} For the two curved mechanisms, OLS-x bias has the sign and size predicted by $g'(p)-(g(1)-g(0))$: in 30 (graph, mechanism, $p$) cases of Experiment 9 the empirical noiseless bias and the parameter-free prediction have correlation $1.00$, mean absolute difference $0.021$ and maximum $0.07$, with every difference within 1.7 Monte-Carlo standard errors (Table~\ref{tab:theory} in the appendix). The bias flips sign: for \texttt{sat} it is $+0.86$ at $p{=}0.2$ and $-0.37$ at $p{=}0.5$, and for \texttt{thresh} it is $-0.83$ at $p{=}0.2$ and $+0.25$ at $p{=}0.5$ (SBM, Exp.~9). With $p\le0.1$ the linear model's RMSE is $1.7$--$2.2$ for \texttt{sat} (GATE $=2$). At $p{=}0.5$ the quadratic $x^2$ is estimated without bias, purely because tangent and secant slopes coincide. The bias does not decrease with data: for \texttt{sat} under Bernoulli designs OLS-x bias is $-0.33,-0.38,-0.41,-0.37,-0.39$ for $n=1$k$,2$k$,4$k$,8$k$,16$k, and its RMSE plateaus at $0.42$ while the oracle model's RMSE keeps falling to $0.22$ (Appendix Fig.~\ref{fig:e4}). With fixed $n{=}3000$, higher degree makes things worse (Fig.~\ref{fig:e6}): OLS-x RMSE (Bernoulli, \texttt{sat}) rises from $0.17$ ($\bar d{=}4$) to $2.9$ ($\bar d{=}128$), exceeding DM ($1.00$) between $\bar d{=}32$ ($0.86$) and $\bar d{=}64$ ($1.38$); the noise term scales as $\sqrt{d/n}$, so the relevant quantity is $d/n$.

\begin{figure*}[t]
\centering
\includegraphics[width=\textwidth]{../figs/fig_e6.pdf}
\caption{Degree sweep on SBM graphs ($n{=}3000$, 15 blocks, 200 replicates). Top: bias of OLS-x from noiseless fits $\pm1.96$ MC s.e.; dotted: Prop.~\ref{prop:bern}; dashed grey: $d\to\infty$ limit. Bottom: RMSE.}\label{fig:e6}
\end{figure*}

\begin{figure*}[t]
\centering
\includegraphics[width=\textwidth]{../figs/fig_e3.pdf}
\caption{Design comparison on SBM graphs when the between-block edge share grows ($n{=}3000$, 200 replicates, 95\% bootstrap CIs). Rows: mechanism; columns: estimator. Relative RMSE (RMSE/|GATE|).}\label{fig:e3}
\end{figure*}

\textbf{R4: Design matters mostly through cluster quality (Fig.~\ref{fig:e3}, Table~\ref{tab:e3}).} Cluster randomisation removes most of the extrapolation: on SBM with 10\% between-block edges, \texttt{sat} OLS-x has bias $+0.01$ (vs.\ $-0.39$ for Bernoulli(0.5)) and $\|w\|_1=2.4$ (vs.\ $6.4$); absolute RMSE (GATE $=2$) falls from 0.55 to 0.22, and DM's from 1.0 to 0.26. Its advantage erodes as clusters get worse: at 60\% between-block edges the cluster design's OLS-x RMSE (0.50) is worse than Bernoulli(0.5) (0.43), and randomised saturation lies between Bernoulli and cluster randomisation for $\le30\%$ between-block edges (OLS-x RMSE $0.40$--$0.49$) and is worst at 60\% ($0.82$); it was never better than plain cluster randomisation here. A small $p$ (0.2) is clearly harmful for regression (bias $+0.88$ for \texttt{sat}), and under Bernoulli(0.2) the H\'ajek estimator found no treated-and-saturated and control-and-isolated units in 98\% of replicates (undefined); for \texttt{thresh} even the \emph{correctly specified} model has RMSE $1.4$--$2.1$ at $p{=}0.2$ because the design never excites the steep region (large $\|w\|$): correct specification is not sufficient. The cluster design retains a variance floor from the few randomisation units (K=15--30): its RMSE does not fall with degree, so it is better than Bernoulli at high degree and worse at low degree (Fig.~\ref{fig:e6}; for OLS-x the crossover is near $\bar d\approx8$--$16$ with 15 blocks).

\begin{table}[t]
\caption{Error anatomy of OLS-x on SBM ($n{=}3000$, 10\% between-block edges): bias, SD of the assignment-induced part $S$, SD of the noise part $N$, RMSE, and realised weight norms (200 replicates).}\label{tab:e3}
\centering\scriptsize
\setlength{\tabcolsep}{2.5pt}
\begin{tabular}{llrrrrrr}
\toprule
Mech. & Design & bias & SD$_S$ & SD$_N$ & RMSE & $\|w\|_1$ & $\|w\|_2$ \\
\midrule
sat & Bernoulli $p{=}.5$ & -0.39 & 0.35 & 0.15 & 0.55 & 6.4 & 0.147 \\
sat & Bernoulli $p{=}.2$ & +0.88 & 0.42 & 0.18 & 0.99 & 8.0 & 0.184 \\
sat & cluster & +0.01 & 0.21 & 0.05 & 0.22 & 2.4 & 0.049 \\
sat & rand. saturation & -0.19 & 0.38 & 0.07 & 0.43 & 3.2 & 0.067 \\
thresh & Bernoulli $p{=}.5$ & +0.25 & 0.38 & 0.13 & 0.46 & 6.4 & 0.147 \\
thresh & Bernoulli $p{=}.2$ & -0.85 & 0.43 & 0.18 & 0.97 & 8.1 & 0.185 \\
thresh & cluster & -0.00 & 0.22 & 0.05 & 0.22 & 2.4 & 0.049 \\
thresh & rand. saturation & +0.21 & 0.40 & 0.07 & 0.46 & 3.2 & 0.068 \\
hetero & Bernoulli $p{=}.5$ & -0.04 & 0.35 & 0.14 & 0.37 & 6.4 & 0.147 \\
hetero & Bernoulli $p{=}.2$ & -0.07 & 0.43 & 0.19 & 0.46 & 8.0 & 0.184 \\
hetero & cluster & +0.00 & 0.23 & 0.05 & 0.24 & 2.4 & 0.049 \\
hetero & rand. saturation & -0.01 & 0.36 & 0.06 & 0.36 & 3.2 & 0.067 \\
\bottomrule
\end{tabular}

\end{table}


\textbf{R4b: Regression tolerates finer clusters than DM does (Appendix Table~\ref{tab:e10}).} When the cluster resolution is increased (more, smaller clusters; higher cut) DM degrades monotonically on ego-Facebook (RMSE $0.34\to0.56$ from $K{=}16$ to $672$, \texttt{sat}), whereas OLS-x is non-monotone but ends lower: $0.34$ ($K{=}16$, cut $0.04$), $0.41$ ($K{=}27$), $0.26$ ($K{=}238$, cut $0.53$); similarly on Stanford3 ($0.60\to0.41$ from $K{=}10$ to $197$) and WS ($0.24\to0.17$, $K{=}27\to156$). It does not help indefinitely: on SBM with 30\% between-block edges OLS-x worsens ($0.23\to0.49$) when Louvain splits the blocks ($K{=}376$, cut $0.78$). The right granularity trades the variance floor from few randomisation units against the cut and differs between DM and regression.

\textbf{R4c: Outcome homophily drives the variance of cluster designs (Appendix Table~\ref{tab:e11}).} Our default community-level baseline effect (sd $0.5$) is one choice. Varying it from 0 to 2 (SBM, \texttt{sat}) changes OLS-x RMSE under cluster randomisation from $0.06$ to $0.84$ (and under Bernoulli from $0.41$ to $1.37$), while DM under Bernoulli stays at $1.00$ (pure bias). With strong community-level baselines, Bernoulli-design OLS-x ($1.37$) is worse than DM ($1.00$), and the accuracy advantage of the cluster design over DM under cluster randomisation shrinks ($0.84$ vs.\ $0.87$). The headline accuracy numbers (R1) therefore depend on how much baseline outcome variation is shared within clusters, a quantity the analyst can estimate from pre-experiment data but we did not exploit.

\textbf{R5: Flexible learners and richer features do not fix extrapolation (Table~\ref{tab:e2}, Appendix Fig.~\ref{fig:e2app}).} Gradient boosting has median relative RMSE 0.23--0.44 under Bernoulli designs (OLS-x: 0.16--0.32), i.e.\ worse than OLS-x on all six graphs; on the synthetic graphs its bias is towards zero ($-0.2$ to $-0.4$ median; tree predictions cannot leave the range of observed exposures) while on the two real graphs it is near zero ($+0.02$) with a large SD, and its SD is $1.2$--$2.6\times$ that of OLS-x; it does not converge with $n$ in the range tried (RMSE $0.65$ at $n{=}16$k for \texttt{sat}, bias $-0.47$). The 10-feature polynomial model \texttt{rich} is the worst regression (median over cells $0.86$ vs.\ $0.21$ for OLS-x): extra features inflate $\|w\|_2$ (mean $0.71$ vs.\ $0.12$) more than they remove bias. Adding $x^2$ and $zx$ (\texttt{quad}) is neutral (median $0.21$ vs.\ $0.21$), and the two-hop feature removes the bias of the \texttt{twohop} mechanism ($+0.07$ vs.\ $-0.47$, ER/Bernoulli) but its variance ($\|w\|_2=0.36$) makes the RMSE higher than OLS-x in all 12 design$\times$graph cells ($1.11$ vs.\ $0.54$ on ER/Bernoulli): bias--variance, not a free fix.

\textbf{R6: Regression only pays off above a spillover threshold (Appendix Fig.~\ref{fig:e5}).} With $\gamma{=}0$ DM has RMSE $0.06$ against $0.37$ for OLS-x (linear mechanism, Bernoulli): the price of the regression is $\approx6\times$ larger error when there is no interference; DM wins at $\gamma{=}0$ and OLS-x at $\gamma{=}0.5$ ($0.34$ vs.\ $0.51$), so the break-even lies in between, near $\gamma\approx0.35$ by linear interpolation (spillover $\approx26\%$ of the GATE; we did not run intermediate values). Under cluster randomisation on the SBM graph DM and OLS-x are comparable for small $\gamma$ ($0.22$ vs.\ $0.22$ at $\gamma{=}0$) and OLS-x is better for $\gamma\ge2$ (e.g.\ $0.24$ vs.\ $0.44$ for \texttt{sat}).

\textbf{R7: Wrong graph.} (Fig.~\ref{fig:e8}.) If the analyst's graph deletes a fraction $r$ of the true edges and adds as many random ones, the Bernoulli-design bias grows roughly like $-r$: OLS-x bias for \texttt{linfrac} is $-0.07,-0.24,-0.50,-0.75$ for $r=0.1,0.25,0.5,0.75$ (RMSE $0.34\to0.78$ for $r$ from 0 to 0.75), approaching DM's. Under cluster randomisation the estimator is robust ($-0.01$ at $r{=}0.25$, $-0.08$ at $r{=}0.75$) because $x$ is dominated by own-cluster treatment. Using the correct feature form does not remove this bias when the graph is wrong (oracle $\approx$ OLS-x for \texttt{linfrac}; for \texttt{sat} oracle RMSE $0.62$ vs.\ $0.73$ at $r{=}0.5$).



\textbf{R8: Uncertainty quantification is the weakest link (Table~\ref{tab:e7}).} For OLS-x on 16 populations, 95\% intervals from the usual i.i.d.\ model SE or heteroskedasticity-robust (HC0/HC3) SEs under-cover even when the model is correct: for \texttt{linfrac} on the SBM graph coverage is $0.76$ (Bernoulli) and $0.48$ (cluster) as the SEs are $\approx0.60$ and $0.32\times$ the true SD, because residuals are positively correlated through neighbouring units and shared community effects. SEs clustered on the design clusters (CR) restore near-nominal coverage when the clusters are meaningful (SBM, correct model: $0.93$, $0.92$) and fail when they are not (ER: $0.89$, $0.82$). When the model is biased (\texttt{sat}, \texttt{peer}; Appendix Table~\ref{tab:e7full}) coverage of the i.i.d.\ intervals falls to $0.34$--$0.61$ and even CR coverage is as low as $0.53$ (ER, \texttt{peer}).

\begin{table}[t]
\caption{Coverage of nominal 95\% intervals for the GATE (OLS-x, $n{=}3000$, 500 replicates) with i.i.d., HC3 and design-cluster-robust (CR) SEs; last two columns: mean SE/true SD of the estimator.}\label{tab:e7}
\centering\scriptsize
\setlength{\tabcolsep}{2pt}
\begin{tabular}{lllrrrrrrr}
\toprule
Graph & Mech. & Des. & bias & RMSE & iid & HC3 & CR & $\overline{\rm SE}_{iid}$/SD & $\overline{\rm SE}_{CR}$/SD \\
\midrule
SBM & linfrac & Bern & +0.01 & 0.36 & 0.76 & 0.76 & 0.93 & 0.59 & 0.97 \\
SBM & linfrac & Clust & +0.00 & 0.22 & 0.48 & 0.48 & 0.92 & 0.32 & 0.93 \\
SBM & sat & Bern & -0.35 & 0.50 & 0.55 & 0.56 & 0.82 & 0.58 & 0.98 \\
SBM & sat & Clust & -0.02 & 0.23 & 0.45 & 0.46 & 0.90 & 0.31 & 0.91 \\
ER & linfrac & Bern & +0.01 & 0.27 & 0.91 & 0.91 & 0.89 & 0.84 & 0.84 \\
ER & linfrac & Clust & +0.01 & 0.32 & 0.79 & 0.79 & 0.82 & 0.64 & 0.69 \\
ER & sat & Bern & -0.37 & 0.47 & 0.61 & 0.63 & 0.59 & 0.80 & 0.81 \\
ER & sat & Clust & -0.35 & 0.46 & 0.56 & 0.57 & 0.60 & 0.68 & 0.74 \\
\bottomrule
\end{tabular}

\end{table}


\textbf{R9: Diagnostics (Fig.~\ref{fig:diag}).} $\|w\|_2$ predicts the noise SD almost perfectly (R2), so a practitioner can compute $\hat\sigma\|w\|_2$ from the realised assignment. Across all configurations with a correct graph, $\|w\|_1$ rank-correlates with $|$bias$|$ of OLS-x for curved responses (Spearman $0.81$ for \texttt{sat}, $0.83$ for \texttt{thresh}, $0.75$ for \texttt{twohop}; 112, 80, 30 configurations) and weakly where the model is (nearly) correct ($0.21$ for \texttt{linfrac}, where the true bias is zero and the correlation reflects only Monte-Carlo noise); it is a risk indicator, not a bias estimate.

\begin{figure*}[t]
\centering
\begin{minipage}[t]{0.58\textwidth}\centering
\includegraphics[width=\linewidth]{../figs/fig_e8.pdf}
\captionof{figure}{RMSE when the analyst's graph is corrupted (fraction of true edges replaced by random ones), SBM, 200 replicates; top: Bernoulli, bottom: cluster randomisation. Oracle uses the true exposure form on the corrupted graph.}\label{fig:e8}
\end{minipage}\hfill
\begin{minipage}[t]{0.38\textwidth}\centering
\includegraphics[width=\linewidth]{../figs/fig_diag.pdf}
\captionof{figure}{Top: empirical SD of the noise term vs.\ $\sigma\|w\|_2$ (all OLS configurations). Bottom: $|$bias$|$ of OLS-x vs.\ $\|w\|_1$ (correct graph).}\label{fig:diag}
\end{minipage}
\end{figure*}
